% ===== latex-recomposed:document =====
% flm: style file mydoc.7d2bb188.sty
\flmRef{ref1} \flmHeading[toc=true]
\section{A}{\flmLabels{\label{ref1}}} \begin{align} x \end{align} \EscVerb[formatcom=\flmFmtVRB{verbtext}]{v}%%
% ===== latex-recomposed:style mydoc.7d2bb188.sty =====
\NeedsTeXFormat{LaTeX2e}
\ProvidesPackage{mydoc.7d2bb188}
% --- packages ---
% flm.latex.pkg.fvextra
\RequirePackage{fvextra}
% --- commands and environments the library defines for the document ---
% flm.latex.def.refs
% \flmRef[<options>]{<label>}: a reference that gives the name of its target
% ("Section 2"), set by \zcref of the zref-clever package.
%
% The options are the ones the document wrote, in FLM's own vocabulary. This
% definition translates each of them by its key -- a value written after the
% key is ignored, as in FLM -- and drops every key it has no translation for:
%
%   noprefix    -> noname                 the number or title alone, with no
%                                         name before it ("2", not
%                                         "Section 2")
%   S           -> S                      the first name capitalized and
%                                         unabbreviated, for a reference that
%                                         opens a sentence
%   capital     -> cap=true,abbrev=false  the names capitalized ("Lemma 1")
%   lowercase   -> cap=false,abbrev=false the names in lowercase ("lemma 1")
%   abbreviated -> abbrev=true            the abbreviated names ("Lem. 1")
%                                         where the reference type declares
%                                         them, the full names otherwise
%   anything else -> dropped              the default wording; a variant of
%                                         your own naming has no counterpart
%                                         in zref-clever, which would stop on
%                                         a key it does not know
%
% Redefine \flmRef to set references differently -- another package, another
% wording, your own variants, no link. Your definition sees the options as the
% document wrote them, in FLM's vocabulary, and maps them as it likes.
%
% The translation uses expl3, LaTeX's programming layer, which is part of the
% format and needs no package; its names read only between \ExplSyntaxOn and
% \ExplSyntaxOff. \keyval_parse:NNn reads the option list as LaTeX reads any
% key-value list: it hands a key written alone to the first macro below and a
% key with its value to the second, which drops the value. \str_case:nnF looks
% the key up in the table and appends what \zcref is to be given, or nothing.
% The macros with @ in their names are internals, not a customization point:
% redefine \flmRef, not its parts.
\ExplSyntaxOn
\def\flm@ref@add#1{\tl_put_right:Nn \flm@ref@options {#1,}}
\def\flm@ref@key#1{\str_case:nnF {#1}
  {{noprefix}{\flm@ref@add{noname}}
   {S}{\flm@ref@add{S}}
   {capital}{\flm@ref@add{cap=true,abbrev=false}}
   {lowercase}{\flm@ref@add{cap=false,abbrev=false}}
   {abbreviated}{\flm@ref@add{abbrev=true}}}
  {}}
\def\flm@ref@pair#1#2{\flm@ref@key{#1}}
\def\flm@ref@zcref#1#2{\zcref[#1]{#2}}
\ProvideDocumentCommand\flmRef{O{} m}{%
  \def\flm@ref@options{}%
  \keyval_parse:NNn \flm@ref@key \flm@ref@pair {#1}%
  \exp_args:No \flm@ref@zcref {\flm@ref@options} {#2}}
\ExplSyntaxOff
% flm.latex.def.label-display-text
% \flmLabelsWithDisplayText{<text>}{<labels>}: makes <labels> refer to this
% place, with <text> as what a reference to them prints. It is the only
% command for that: a heading the class did not number uses it, and so does
% anything else that has a name but no number.
%
% It places no anchor of its own. Under hyperref a sectioning command, starred
% or not, already places one at the title; a \label after it points there.
% \refstepcounter would place a second anchor here, below the title, and a
% link would lead below the heading. So the counter is stepped with
% the plain \stepcounter and the two things this needs are set here:
% \@currentlabel, which is what a \label records, and \@currentcounter,
% which is what zref-clever reads to give the label its reference type. The
% group keeps every change local.
%
% <text> is expanded to make \@currentlabel, and \label does not survive that
% expansion; so a \label inside <text> (for example a footnote's) is dropped
% from the label's text only. At its own place in the document it still makes
% its own label.
%
% A caller with no sectioning command before it -- for example a defined term
% in running text -- writes \phantomsection at its own place, so that a link
% has an anchor of its own to lead to.
%
% Redefine \flmLabelsWithDisplayText to change how such labels are made.
\ifcsname c@flm@internal@refcnt\endcsname\else\newcounter{flm@internal@refcnt}\fi
\ProvideDocumentCommand\flmLabelsWithDisplayText{m m}{%
  \begingroup
    \stepcounter{flm@internal@refcnt}%
    {\let\label\@gobble\protected@xdef\flm@lwdt@text{#1}}%
    \let\@currentlabel\flm@lwdt@text
    \def\@currentcounter{flm@internal@refcnt}%
    #2%
  \endgroup}
% flm.latex.def.zref-clever-type
% The zref-clever reference type of the counter \flmLabelsWithDisplayText
% steps, with every name prefix empty, so that a reference to such a label
% prints its text alone ("Introduction", not "Section Introduction").
% Every other label keeps the name its own type gives it. Redefine these
% settings, or the type of the counter, to put something before the text.
\zcsetup{countertype={flm@internal@refcnt=flm--internal--ref}}
\zcRefTypeSetup{flm--internal--ref}{Name-sg = , name-sg = , Name-pl = , name-pl = ,}
% flm.latex.def.labels
% \flmLabels{<\label commands>}: the labels of one place in the document.
% In text it runs them as written. On a line of a display-math
% environment, where amsmath keeps the last label of a line and loses the
% ones before it, it runs the first and records every further one as an alias:
% a line per alias is written to the .aux file after the last page and read on
% the next run, making the alias refer to what the first label refers to
% (for \ref and for \zcref alike). Redefine \flmLabels to change what a
% group of labels does; \def\flmLabels#1{#1} restores plain LaTeX behavior,
% amsmath's "Multiple \label's" message included.
\ProvideDocumentCommand\flmLabels{m}{%
  \ifmmode\expandafter\flm@labels@math\else\expandafter\@firstofone\fi{#1}}
% Reading the group: the first \label as written, the rest as aliases of it.
\def\flm@labels@math#1{\def\flm@lbl@first{}\flm@lbl@scan#1\flm@lbl@stop}
\def\flm@lbl@scan{\@ifnextchar\label\flm@lbl@take\flm@lbl@rest}
\def\flm@lbl@take\label#1{%
  \ifx\flm@lbl@first\@empty
    \def\flm@lbl@first{#1}%
    \label{#1}%
  \else
    \expandafter\flm@recordalias\expandafter{\flm@lbl@first}{#1}%
  \fi
  \flm@lbl@scan}
\def\flm@lbl@rest#1\flm@lbl@stop{#1}
% Collecting the aliases. A name is recorded once however often the line is
% typeset (amsmath measures a line before setting it).
\def\flm@aliaslist{}
\def\flm@recordalias#1#2{%
  \ifcsname flm@alias@#2\endcsname\else
    \expandafter\gdef\csname flm@alias@#2\endcsname{}%
    \g@addto@macro\flm@aliaslist{\flm@alias@entry{#2}{#1}}%
  \fi}
% Writing them, from the hook that runs after the last page has been shipped
% out: every \newlabel of that page has reached the .aux file by then, so the
% alias lines come after them. \AtEndDocument would run too early.
\def\flm@alias@write#1#2{\immediate\write\@auxout{\string\flm@aliaslabel{#1}{#2}}}
\AddToHook{enddocument/afterlastpage}{\let\flm@alias@entry\flm@alias@write\flm@aliaslist}
% Reading them back: the alias refers to what the original refers to. The
% assignments are global because the .aux file is read inside a group.
\def\flm@aliaslabel#1#2{%
  \ifcsname r@#2\endcsname
    \global\expandafter\let\csname r@#1\expandafter\endcsname\csname r@#2\endcsname
  \fi
  \ifcsname Z@R@#2\endcsname
    \global\expandafter\let\csname Z@R@#1\expandafter\endcsname\csname Z@R@#2\endcsname
  \fi}
% flm.latex.def.heading
% \flmHeading[<keys>]<sectioning command><star>{<title>}{<labels>}: a heading
% and the labels that refer to it, written on two lines for readability, as in
%   \flmHeading[toc=true]
%   \section*{Title}{\flmLabels{\label{ref1}}}
% It runs the sectioning command as written, then makes the labels refer to
% the heading: by its number when the class numbered it, by its title when it
% did not (a starred heading, or a level below secnumdepth). It finds out
% which without patching anything: a class numbers a heading through
% \refstepcounter, which records the counter's name in \@currentcounter.
%
% The optional argument is a key-value list. It is written here, on the hook,
% and never on the sectioning command: a starred sectioning command takes no
% optional argument, and giving one to \section* is not an error, it is
% silently wrong -- \section*[Short]{Long} takes one token as its heading,
% so the opening bracket alone becomes the heading and Short, the closing
% bracket and Long all run together into the following paragraph. Two keys
% are defined:
%
%   toc=true|false     whether the heading is listed in the table of contents.
%                      Absent, it is listed as LaTeX lists it. After a starred
%                      sectioning command with toc=true, which LaTeX never
%                      lists, \addcontentsline{toc}{<level>}{<entry>} adds it:
%                      <level> is the command's own name (section for
%                      \section), <entry> the alternative title if alt was
%                      written, else the title; under hyperref the starred
%                      command has already placed the anchor the line links
%                      to. On an unstarred heading toc=false is not honored:
%                      the heading stays listed and a warning says so.
%
%   alt={Short title}  an alternative title: what goes into the table of
%                      contents, the running head and the PDF bookmark, while
%                      the page keeps the full title. It is handed to the
%                      sectioning command as its own optional argument, so it
%                      does what \section[Short title]{Full title} does, and
%                      nothing else: a reference to an unnumbered heading still
%                      prints the full title. A starred heading does not hand
%                      it to the sectioning command; where it is written
%                      toc=true, alt is its line in the table of contents.
%
% What counts is whether the key was written, not what it holds. alt={} is an
% empty alternative title and is handed over as \section[]{Full title}, which
% is how LaTeX asks for an empty table-of-contents entry; no optional argument
% at all, and an empty one ([]), hand nothing over and leave the sectioning
% command exactly as it was. Writing alt with no value is LaTeX's own error,
% as is a key that is not defined here, and a value of toc other than true or
% false. The values and the record that alt was written are cleared at every
% call, so one heading's keys never reach the next. The recomposed LaTeX writes
% toc on every heading and alt where the heading was given an alternative
% title; both keys are also there for a document written by hand and for a
% user who redefines the hook.
%
% Redefine \flmHeading to change what a heading does. The block below uses
% expl3, LaTeX's programming layer, which is part of the format and needs no
% package; its names read only between \ExplSyntaxOn and \ExplSyntaxOff.
% \keys_define:nn declares the keys and \keys_set:nn reads a list of them;
% \bool_if:NTF asks whether the key was written. The three declarations are
% guarded on their own names, as every declaration in this file is, so that two
% of these style files can be loaded in one LaTeX run. The macros with @ in
% their names are internals, not a customization point: redefine \flmHeading,
% not its parts.
\ExplSyntaxOn
\ifcsname flm@heading@alt\endcsname\else\tl_new:N \flm@heading@alt\fi
\ifcsname flm@heading@given\endcsname\else\bool_new:N \flm@heading@given\fi
\ifcsname flm@heading@toc\endcsname\else\tl_new:N \flm@heading@toc\fi
\keys_define:nn {flm/heading}
  {alt .code:n = {\tl_set:Nn \flm@heading@alt {#1}
                  \bool_set_true:N \flm@heading@given} ,
   alt .value_required:n = true ,
   toc .choices:nn = {true, false} {\tl_set_eq:NN \flm@heading@toc \l_keys_choice_tl} ,
   toc .value_required:n = true}
\def\flm@heading@withalt#1#2#3{#2[#1]{#3}}
% The contents line of a starred heading: <level> then <entry>.
\def\flm@heading@addtoc#1#2{\addcontentsline{toc}{#1}{#2}}
% toc=false on an unstarred heading: the heading stays listed.
\def\flm@heading@notoc#1{\PackageWarning{flm}{toc=false~on~the~unstarred~
  \string#1~is~not~honored~by~the~default~\string\flmHeading;~the~heading~
  stays~in~the~table~of~contents}}
\ProvideDocumentCommand\flmHeading{O{} m s m m}{%
  \tl_clear:N \flm@heading@alt
  \bool_set_false:N \flm@heading@given
  \tl_clear:N \flm@heading@toc
  \keys_set:nn {flm/heading} {#1}%
  \def\@currentcounter{}%
  \IfBooleanTF{#3}
    {#2*{#4}%
     \str_if_eq:VnT \flm@heading@toc {true}
       {\edef\flm@heading@level{\cs_to_str:N #2}%
        \bool_if:NTF \flm@heading@given
          {\exp_args:NVV \flm@heading@addtoc \flm@heading@level \flm@heading@alt}
          {\exp_args:NV \flm@heading@addtoc \flm@heading@level {#4}}}}
    {\str_if_eq:VnT \flm@heading@toc {false} {\flm@heading@notoc{#2}}%
     \bool_if:NTF \flm@heading@given
       {\exp_args:NV \flm@heading@withalt \flm@heading@alt #2{#4}}
       {#2{#4}}}%
  \ifx\@currentcounter\@empty
    \flmLabelsWithDisplayText{#4}{#5}%
  \else
    #5%
  \fi}
\ExplSyntaxOff
% flm.latex.def.verbatim
% Verbatim content. Inline content is written as
%   \EscVerb[formatcom=\flmFmtVRB{<kind>}]{<escaped content>}
% and block content as the flmVerbatimCode environment. \EscVerb comes with
% fvextra 1.5 (2022) and later, which this file requires.
% \flmFmtVRB{<kind>} runs the hook \flmFmtVRB<kind>; redefine
% \flmFmtVRBverbtext, \flmFmtVRBverbcode or \flmFmtVRBverba to change how
% that kind of verbatim content is set.
\ProvideDocumentCommand\flmFmtVRBverbtext{}{}
\ProvideDocumentCommand\flmFmtVRBverbcode{}{\ttfamily}
\ProvideDocumentCommand\flmFmtVRBverba{}{\itshape}
\ProvideDocumentCommand\flmFmtVRB{m}{\csname flmFmtVRB#1\endcsname}
% \flmVerbatimCodeSetup{<lang>} runs at the start of a flmVerbatimCode
% environment with its language tag (empty when none was given); redefine
% it to attach syntax highlighting. It runs before the Verbatim environment
% opens, so it must not open a group or an environment of its own.
\ProvideDocumentCommand\flmVerbatimCodeSetup{m}{}
% flmVerbatimCode[<lang>]: a block of code, set with the verbcode hook. The
% recomposed LaTeX always writes the brackets, [<lang>] or [], because LaTeX's
% lookahead for an absent optional argument swallows the line end that the
% verbatim content must begin after. Define the environment before loading
% this file to replace it.
\ProvideDocumentEnvironment{flmVerbatimCode}{O{}}{%
  \flmVerbatimCodeSetup{#1}%
  \VerbatimEnvironment
  \begin{Verbatim}[formatcom=\flmFmtVRB{verbcode}]}%
  {\end{Verbatim}}
% flmVerbatimText[<lang>]: a block of plain verbatim text, set with the
% verbtext hook. It takes the same optional argument as flmVerbatimCode and
% ignores it, so that one command (\flmVerbatimBlock below), which always
% writes the brackets, serves either environment.
\ProvideDocumentEnvironment{flmVerbatimText}{O{}}{%
  \VerbatimEnvironment
  \begin{Verbatim}[formatcom=\flmFmtVRB{verbtext}]}%
  {\end{Verbatim}}
% \flmVerbatimBlock[<options>]{<environment>}{<lines>}: one block of verbatim
% content written as a single macro call, for the places a block cannot be
% written out in full -- inside an argument of a call, inside the body a
% body-grabbing environment hands over, inside a definition's own body -- where
% LaTeX has read the characters long before anything can say they are verbatim.
% The three arguments:
%   <options>     the block environment's own optional argument, always
%                 written, [<lang>] or [];
%   <environment> the environment's name, flmVerbatimCode for a code block;
%   <lines>       one brace group per line of the block, and inside each line
%                 one brace group per tab-separated segment, every segment
%                 written in the encoding \EscVerb reads (a backslash before
%                 every printable ASCII character that is not a letter or a
%                 digit). The line and segment structure is carried by braces
%                 rather than by characters because a line feed written into an
%                 auxiliary file comes back from it as a space, and a tab
%                 inside an argument is a space before anything can look at it;
%                 brace groups survive both.
% The command un-escapes every segment with fvextra's own un-escaper, joins the
% segments of a line with a tab and the lines with a line feed, and hands
% \begin{<environment>}[<options>] ... \end{<environment>} to \scantokens,
% which offers that text to TeX as if it were a file: the environment is
% entered for real, and fancyvrb reads its lines the way it reads them out of
% any file. This is the only place this file uses \scantokens, and it is the
% whole mechanism -- nothing is saved anywhere, no name is made up, nothing
% outlives the call. The \begin and the \end are written as plain characters
% (\string\begin\string{...) because \scantokens puts a space after every
% control word, and fancyvrb looks for the exact characters \end{ on a line:
% written as control words they would read \end {...} and the environment
% would never close.
% The call leaves the mode it found. A block is vertical material, so the
% environment leaves TeX in vertical mode regardless of what came before; a
% block that ends a section title would then leave hyperref's link open in
% vertical mode, which pdfTeX refuses on the run that reads the table of
% contents back. Called
% where a paragraph was under way, the command therefore begins one again
% (\leavevmode); called where TeX was already in vertical mode, it does not,
% because starting a paragraph there would add an empty line the same block
% written out in full does not add.
\begingroup
\catcode`\^^J=12 \catcode`\^^I=12
\gdef\flm@vrb@nl{^^J}\gdef\flm@vrb@tab{^^I}
\endgroup
\let\flm@vrb@stop\relax
\let\flm@vrb@leave\relax
\ProvideDocumentCommand\flmVerbatimBlock{O{} m m}{%
  \ifhmode\let\flm@vrb@leave\leavevmode\else\let\flm@vrb@leave\relax\fi
  \begingroup
    \gdef\flm@vrb@acc{}\def\flm@vrb@lsep{}%
    \flm@vrb@lines#3\flm@vrb@stop
    \xdef\flm@vrb@src{%
      \string\begin\string{#2\string}[#1]\flm@vrb@nl
      \flm@vrb@acc\flm@vrb@nl\string\end\string{#2\string}}%
  \endgroup
  \edef\flm@vrb@restorenl{\newlinechar=\the\newlinechar\relax}%
  \newlinechar=`\^^J\relax
  \expandafter\scantokens\expandafter{\flm@vrb@src}%
  \flm@vrb@restorenl
  \ifvmode\flm@vrb@leave\fi}
% The walk over the lines and their segments: one line group at a time, one
% segment at a time, the segments of a line joined with a tab and the lines
% with a line feed. \FVExtraDetokenizeREscVArg is fvextra's own un-escaper,
% the macro \EscVerb itself calls, so that inline verbatim content and a block
% share one escaping; it is also the only safe route for text that is not
% ASCII.
\def\flm@vrb@lines{\@ifnextchar\bgroup\flm@vrb@line\flm@vrb@linesdone}
\def\flm@vrb@line#1{%
  \xdef\flm@vrb@acc{\flm@vrb@acc\flm@vrb@lsep}\let\flm@vrb@lsep\flm@vrb@nl
  \def\flm@vrb@ssep{}\flm@vrb@segs#1\flm@vrb@stop\flm@vrb@lines}
\def\flm@vrb@linesdone#1\flm@vrb@stop{}
\def\flm@vrb@segs{\@ifnextchar\bgroup\flm@vrb@seg\flm@vrb@segsdone}
\def\flm@vrb@seg#1{\FVExtraDetokenizeREscVArg{\flm@vrb@addseg}{#1}\flm@vrb@segs}
\def\flm@vrb@segsdone#1\flm@vrb@stop{}
\def\flm@vrb@addseg#1{%
  \xdef\flm@vrb@acc{\flm@vrb@acc\flm@vrb@ssep#1}\let\flm@vrb@ssep\flm@vrb@tab}
% What a PDF bookmark holds where a heading contains verbatim content. A
% bookmark is plain text: it can carry neither the setting of inline verbatim
% content nor a block at all, so hyperref is told what to put there instead.
% Inline content gives its text -- without this, fvextra disables \EscVerb
% without removing its optional argument and the option list appears in the
% bookmark -- and a block gives its content run together, with no line breaks
% and no separator between the lines. The two definitions are in a hook that
% runs when the document begins and does nothing where hyperref is absent, so
% this file neither requires hyperref nor cares whether it is loaded before or
% after it.
\AtBeginDocument{\ifdefined\pdfstringdefDisableCommands
  \pdfstringdefDisableCommands{%
    \def\EscVerb[#1]#2{#2}%
    \def\flmVerbatimBlock[#1]#2#3{#3}%
  }\fi}
\endinput
